\section{Uniform control of returning geodesics}
\label{sec:turning}

The only passages still requiring an estimate are geodesics that enter the
neck and return to the same boundary.  The boundary term in the index form
controls one component of the endpoint change in the torus part of the
variation field.  We show
that this control compensates for the tangential term in
\eqref{eq:turning-index-reduction}, uniformly over all such passages.

\begin{proposition}\label{prop:turning-index}
For all sufficiently small positive $\delta$, every returning geodesic
passage of the constructed neck has nonnegative perpendicular index form
$I^\perp(V)$ for every smooth field $V$ along the closed passage, with no
restriction on its endpoint values.  A single choice of $\delta$ works for
both boundary components and all unit-speed returning geodesics.
\end{proposition}

We first record the estimates on the duration of a passage.  Use the
coordinates of the boundary through which it enters, so that this boundary
is $t=0$ and
\[
 P(0)=\operatorname{diag}(1,b),\qquad
 P'(0)=\kappa ee^\top,\qquad
 e=(1,0)^\top,\qquad \kappa=\delta a(0)>0.
\]
By Lemma~\ref{lem:neck-passages}, the passage occupies $[0,L]$ with
$L>0$, has one turn at $s=L/2$, and satisfies
\[
 t(0)=t(L)=0,\qquad t(L/2)=t_*\in(0,1),\qquad
 v(0)=-v(L)=v_0>0,
\]
where $v=\dot t$, $p=Q(t)\dot x$ is constant, and
$F(t)=p^\top P(t)p$.  In particular,
\begin{equation}\label{eq:turning-basic}
 0<F'(t)\le\kappa p_1^2\quad(0\le t\le t_*),\qquad
 0<v_0^2=F(t_*)-F(0)\le\kappa p_1^2.
\end{equation}
Thus $p_1\ne0$.  Here and throughout this section, primes on $P$ and $F$
denote differentiation in $t$, and dots denote differentiation in $s$.

\begin{lemma}\label{lem:turning-time}
For every returning passage,
\begin{equation}\label{eq:turning-chord}
 \int_0^L v(s)^2\,ds\le \frac{Lv_0^2}{3}.
\end{equation}
If $t_*>1/4$, then
\begin{equation}\label{eq:turning-time-margin}
 4v_0\le\beta\kappa p_1^2L,
 \qquad \beta=\frac{1+\exp(-9/26)}{2}<1.
\end{equation}
\end{lemma}

\begin{proof}
On $[0,L/2]$, $v\ge0$ and
\[
 \ddot v=-\tfrac12F''(t)v\ge0.
\]
Consequently the graph of $v$ lies below the chord joining its endpoint
values $v_0$ and $0$:
\[
 0\le v(s)\le v_0(1-2s/L),\qquad 0\le s\le L/2.
\]
Squaring and integrating on this half-interval, and using
$v(L-s)=-v(s)$, proves \eqref{eq:turning-chord}.

On the same half-interval, $\ddot t=-F'(t)/2<0$, so concavity gives
$t(s)\ge 2st_*/L$.  If $t_*>1/4$, symmetry therefore implies
$t(s)\ge1/8$ throughout $[L/4,3L/4]$.  The end formula for $P$ gives
$F'(1/8)=\delta a(1/8)p_1^2$, and concavity of $F$ then gives
\[
 F'(t(s))\le
 \begin{cases}
 \delta a(1/8)p_1^2,&L/4\le s\le3L/4,\\
 \kappa p_1^2,&0\le s\le L.
 \end{cases}
\]
Since $\dot v=-F'(t)/2$ and $v(L)=-v_0$,
\[
 4v_0=\int_0^L F'(t(s))\,ds
 \le \frac{\delta Lp_1^2}{2}\bigl(a(0)+a(1/8)\bigr).
\]
Finally $a(1/8)/a(0)=\exp(-9/26)$, proving
\eqref{eq:turning-time-margin}.
\end{proof}

We now compare the change in the torus component of a field with its
coordinate energy.  Apply Lemma~\ref{lem:turning-gauge} and write the
resulting field as $V=(w,z)$, so that $w(0)=w(L)=0$ and
$\langle V,T\rangle$ is affine.  Put
\[
 X=z(L)-z(0),\qquad \mu=\frac{\kappa v_0}{4}>0,\qquad
 \xi=\dot z-P'(t)p\,w.
\]
The coordinate energy and the index reduction are
\begin{align}
 E&=\int_0^L\left(|\xi|_Q^2+\dot w^2
                         -\tfrac12F''(t)w^2\right)ds,\notag\\
 I^\perp(V)&\ge E+\mu(e^\top X)^2
                         -\frac{(p^\top X)^2}{L}.
 \label{eq:turning-target}
\end{align}
Thus our target is
\[
 |p^\top X|^2\le L\bigl(E+\mu(e^\top X)^2\bigr).
\]
The boundary penalty controls only one component of $X$, so its
interaction with $p$ matters.

For any constant covector $q\in\mathbb R^2$, integrating
$\dot z=\xi+P'p\,w$ gives
\[
 q^\top X=\int_0^L q^\top\xi\,ds
              +\int_0^L q^\top P'(t)p\,w\,ds.
\]
This identity separates the two terms already present in $E$.
For the second term, equip $\mathcal D=H_0^1(0,L)$ with the inner product
\[
 (w_1,w_2)_{\mathcal D}
 =\int_0^L\left(\dot w_1\dot w_2
                   -\tfrac12F''(t)w_1w_2\right)ds,
\]
and define
\[
 \Phi_q(w)=\int_0^L q^\top P'(t)p\,w\,ds.
\]
Because $-F''\ge0$ and the endpoints vanish, this inner product is
positive definite.  On the fixed finite interval it defines a norm
equivalent to the usual $H_0^1$ norm, by boundedness of $F''(t(s))$
and the Poincar\'e inequality.  No equivalence constant uniform in
$L$ is needed.  In particular, $\Phi_q$ is continuous.

Regard the right side of the functional identity as a functional of
arbitrary pairs $(\xi,w)$ in
$L^2([0,L];\mathbb R^2,Q)\oplus\mathcal D$; the gauged fields provide
particular pairs in this space.  Its squared norm is $S(q,q)$, where
\begin{equation}\label{eq:turning-S}
 S(q,r)=\int_0^L q^\top P(t)r\,ds
                  +(\Phi_q,\Phi_r)_{\mathcal D^*}
\end{equation}
and the dual inner product is that of the Riesz representatives.
Indeed, the first summand of the functional has Riesz representative
$Pq$ in the $Q$ inner product.  The symmetric bilinear form $S$ is
positive definite because its first term is positive definite.

\begin{lemma}\label{lem:turning-duality}
For every constant covector $q$ and every gauged field as above,
\begin{equation}\label{eq:turning-dual-cs}
 |q^\top X|^2\le S(q,q)E.
\end{equation}
For each real $\alpha$, the boundary penalty gives the further bound
\begin{equation}\label{eq:turning-boundary-cs}
 |p^\top X|^2\le
 \left(S(p-\alpha e,p-\alpha e)+\frac{\alpha^2}{\mu}\right)
 \left(E+\mu(e^\top X)^2\right).
\end{equation}
The smallest first factor in this inequality is
\begin{equation}\label{eq:turning-minimum}
 S(p,p)-\frac{S(p,e)^2}{S(e,e)+\mu^{-1}}.
\end{equation}
\end{lemma}

\begin{proof}
The squared norm of $(\xi,w)$ in the Hilbert direct sum above is $E$.
The functional identity and Cauchy--Schwarz therefore give
\eqref{eq:turning-dual-cs}.  To use the boundary term, split
\[
 p^\top X=(p-\alpha e)^\top X+\alpha e^\top X.
\]
Apply \eqref{eq:turning-dual-cs} to the first term, then apply
Cauchy--Schwarz in $\mathbb R^2$ to the two vectors
\[
 (\sqrt E,\sqrt\mu\,|e^\top X|),\qquad
 (\sqrt{S(p-\alpha e,p-\alpha e)},|\alpha|/\sqrt\mu).
\]
This proves \eqref{eq:turning-boundary-cs}, including when either
energy term vanishes.  Its first factor is a quadratic in $\alpha$
with positive leading coefficient $S(e,e)+\mu^{-1}$.  Completing the
square gives \eqref{eq:turning-minimum}, attained at
$\alpha=S(p,e)/(S(e,e)+\mu^{-1})$.
\end{proof}

It now suffices to bound \eqref{eq:turning-minimum} by $L$.
The next lemma allows an excess of $Lv_0^2$ in $S(p,p)$; we will show
that the subtracted boundary correction recovers at least that amount.
For this purpose $S(e,e)$ needs only a bound proportional to $L$,
whereas $S(p,e)$ must remain close to $Lp_1$ with an error proportional
to $|p_1|$.  An error controlled merely by $|p|$ would not suffice when
$p_1$ is small.  The timing margin in
\eqref{eq:turning-time-margin} will then make the recovery uniform over
all returns that leave the endpoint collar.  All constants below
depend only on the constructed cometric family and its fixed cutoff;
they are independent of $p,L,v_0,t_*$ and the boundary end.

\begin{lemma}\label{lem:turning-S-bounds}
There are constants $C_e,C_m>0$, uniform for sufficiently small
$\delta>0$, such that every returning passage satisfies
\begin{align}
 S(p,p)&\le L(1+v_0^2),\label{eq:turning-Spp}\\
 S(e,e)&\le C_e L,\label{eq:turning-See}\\
 |S(p,e)-Lp_1|&\le C_m\delta L|p_1|.\label{eq:turning-Spe}
\end{align}
\end{lemma}

\begin{proof}
For $w\in\mathcal D$, integration by parts and $F'(t)=-2\dot v$ give
\[
 \Phi_p(w)=-2\int_0^L\dot v w\,ds
          =2\int_0^L v\dot w\,ds.
\]
Thus $\|\Phi_p\|_{\mathcal D^*}^2\le4\int_0^L v^2\,ds$.
Since $F(t)=1-v^2$, Lemma~\ref{lem:turning-time} yields
\[
 S(p,p)\le L+3\int_0^L v^2\,ds\le L(1+v_0^2).
\]

Choose a constant $C_0\ge1$ such that the cometric estimates give
\[
 \|P\|\le C_0,\qquad \|P'\|\le C_0\delta,\qquad
 |P'p|^2\le C_0\delta(-F'').
\]
By the last inequality and Cauchy--Schwarz with respect to $ds$,
\[
 |\Phi_e(w)|
 \le\sqrt{C_0\delta}\int_0^L\sqrt{-F''(t)}\,|w|\,ds
 \le\sqrt{2C_0\delta L}\,\|w\|_{\mathcal D}.
\]
This argument does not divide by $-F''$ and also applies on its zero set.
For $\delta\le1$, \eqref{eq:turning-See} follows with $C_e=3C_0$.
The mixed term contributed by the dual inner product satisfies
\begin{align*}
 |(\Phi_p,\Phi_e)_{\mathcal D^*}|
 &\le 2\sqrt{Lv_0^2/3}\sqrt{2C_0\delta L}\\
 &\le 2\sqrt{2C_0a(0)/3}\,\delta L|p_1|,
\end{align*}
where \eqref{eq:turning-basic} was used in the last line.

To estimate the other mixed term, fix $t\in[0,t_*]$ and put
$D(t)=P'(0)-P'(t)$.  Concavity of $P$ implies $D(t)\succeq0$, and
\[
 \|D(t)\|\le2C_0\delta,\qquad
 0\le p^\top D(t)p=\kappa p_1^2-F'(t)\le\kappa p_1^2.
\]
Cauchy--Schwarz for the positive semidefinite form $D(t)$ gives
\[
 |e^\top D(t)p|^2
 \le(e^\top D(t)e)(p^\top D(t)p)
 \le2C_0a(0)\delta^2p_1^2.
\]
This remains valid if $D(t)$ is singular.  Since
$e^\top P'(0)p=\kappa p_1$, it follows that
\[
 |e^\top P'(t)p|
 \le\bigl(a(0)+\sqrt{2C_0a(0)}\bigr)\delta|p_1|.
\]
Integrating from $0$ to $t\le t_*<1$ and using $e^\top P(0)p=p_1$
gives the same bound for $|e^\top P(t)p-p_1|$.
Its integral in $s$, together with the dual mixed bound, proves
\eqref{eq:turning-Spe}; for example, one can take
\[
 C_m=a(0)+\sqrt{2C_0a(0)}+2\sqrt{2C_0a(0)/3}.
\]
\end{proof}

\begin{proof}[Proof of Proposition~\ref{prop:turning-index}]
If $t_*\le1/4$, the entire passage lies in the end region of nonpositive
sectional curvature, and its perpendicular index integrand is
nonnegative.  Assume therefore that $t_*>1/4$, and apply the unrestricted
gauge of Lemma~\ref{lem:turning-gauge}.

By Lemma~\ref{lem:turning-duality}, it suffices to show that
the boundary correction in \eqref{eq:turning-minimum} is at least
$Lv_0^2$, since \eqref{eq:turning-Spp} bounds the first term by
$L(1+v_0^2)$.

Because $p_1\ne0$, define $r=v_0^2/p_1^2>0$.  By
\eqref{eq:turning-basic}, $r\le a(0)\delta$, while
\eqref{eq:turning-time-margin} and $\mu=\kappa v_0/4$ give
$L\mu\ge r/\beta$.  Suppose $C_m\delta\le1/2$.
Using Lemma~\ref{lem:turning-S-bounds} and the fact that
$x\mapsto x/(1+C_e x)$ is increasing for $x\ge0$, we obtain
\begin{align}
 \frac{S(p,e)^2}{S(e,e)+\mu^{-1}}
 &\ge Lp_1^2(1-C_m\delta)^2\frac{L\mu}{1+C_eL\mu}\notag\\
 &\ge Lv_0^2\frac{(1-C_m\delta)^2}{\beta+C_e r}\notag\\
 &\ge Lv_0^2\frac{(1-C_m\delta)^2}
                          {\beta+C_ea(0)\delta}.
 \label{eq:turning-gain}
\end{align}
In addition to the cometric construction's upper bound on $\delta$,
choose
\begin{equation}\label{eq:turning-delta-choice}
 0<\delta\le
 \min\left\{1,\frac1{2C_m},
                \frac{1-\beta}{2C_m+C_ea(0)}\right\}.
\end{equation}
Then
\[
 (1-C_m\delta)^2\ge1-2C_m\delta
                 \ge\beta+C_ea(0)\delta,
\]
so \eqref{eq:turning-gain} is at least $Lv_0^2$.
Together with \eqref{eq:turning-Spp}, this makes
\eqref{eq:turning-minimum} at most $L$.  Equations
\eqref{eq:turning-boundary-cs} and \eqref{eq:turning-target} now give
$I^\perp(V)\ge0$.

Every denominator used here is positive for an actual returning passage:
$L>0$, $v_0>0$, $p_1\ne0$, and $\mu>0$.  The parameter bound
\eqref{eq:turning-delta-choice} contains none of these quantities.
It therefore applies to arbitrarily shallow incidence and arbitrarily
long excursions without a limiting argument.  The cometric bounds are
uniform in the two end frames, so the same choice applies to returns
through either boundary.  Tangential gauge invariance transfers the
conclusion to the original field with its unrestricted endpoint values.
\end{proof}

\begin{proposition}[All neck passages]\label{prop:passage-index}
Choose $\delta>0$ to satisfy Lemma~\ref{lem:cometric-bounds} and
\eqref{eq:turning-delta-choice}, and then choose the boundary periods as in
Section~\ref{sec:cometric}. Let $\gamma$ be any complete unit-speed geodesic
of the resulting closed manifold. For every connected component $J$ of
its visit to the open neck and every smooth compactly supported field
$V$ along $\gamma$, one has
\[
 I_J^\perp(V)\ge0,\qquad I_J(V)\ge0.
\]
\end{proposition}

\begin{proof}
Lemma~\ref{lem:neck-passages} lists every possible passage, including
constant-height geodesics, infinite visits, and all finite returns.
Lemma~\ref{lem:elementary-passages} handles the full-line and monotone
cases, as well as returns contained in an endpoint collar.
Proposition~\ref{prop:turning-index} handles all remaining returns.
These conclusions apply to the restriction of the original field,
with arbitrary values at finite endpoints and zero values sufficiently
far out at infinite endpoints. They use the same choice of $\delta$
in both endpoint frames and do not depend on the boundary periods.
Finally, $I_J(V)=I_J^\perp(V)+\int_J|\dot q_V|^2\,ds\ge0$.
\end{proof}
